% ============================================================================
% Chapter 25: Formatting in MS Word
% ============================================================================
\chapter{Formatting in MS Word}
\label{ch:word-format}

\begin{objectives}
  \item Apply font formatting: bold, italic, underline, font size, font colour
  \item Change text alignment: left, centre, right, justify
  \item Adjust line spacing and paragraph spacing
  \item Use bullet points and numbered lists
\end{objectives}

\bytetip[tip]{Formatting is what turns a boring block of text into a beautiful, readable document. Think of it like decorating a room --- same space, totally different vibe!}

\section{Font Formatting}
\label{sec:font-format}
\index{MS Word!formatting}

% --- Figure 25.2: Font Formatting Examples ---
\begin{figure}[H]
\centering
\begin{tikzpicture}[
  fmtcard/.style={draw=#1, fill=#1!10, rounded corners=6pt, minimum width=2.8cm, minimum height=1.8cm, align=center, font=\small\sffamily, line width=1pt}
]
  \node[fmtcard=ModuleBlue] at (0,0) {Computer\\\footnotesize Normal};
  \node[fmtcard=FunCoral] at (3.3,0) {\textbf{Computer}\\\footnotesize\key{Ctrl}+\key{B}};
  \node[fmtcard=FunPurple] at (6.6,0) {\textit{Computer}\\\footnotesize\key{Ctrl}+\key{I}};
  \node[fmtcard=LabGreen] at (9.9,0) {\underline{Computer}\\\footnotesize\key{Ctrl}+\key{U}};
  \node[fmtcard=NoteOrange] at (13.2,0) {\sout{Computer}\\\footnotesize Strikethrough};
\end{tikzpicture}
\caption{Different font formatting styles applied to the same word}
\label{fig:font-formats}
\end{figure}

\begin{shortcutbox}
\begin{tabular}{ll}
  \key{Ctrl}+\key{B} & Bold \\
  \key{Ctrl}+\key{I} & Italic \\
  \key{Ctrl}+\key{U} & Underline \\
  \key{Ctrl}+\key{Shift}+\key{>} & Increase font size \\
  \key{Ctrl}+\key{Shift}+\key{<} & Decrease font size \\
\end{tabular}
\end{shortcutbox}

\section{Text Alignment}
\label{sec:alignment}
\index{MS Word!alignment}

% --- Figure 25.1: Text Alignment Visual Demo ---
\begin{figure}[H]
\centering
\begin{tikzpicture}
  % Left aligned
  \node[draw=ModuleBlue, fill=ModuleBlue!5, rounded corners=4pt, minimum width=3cm, minimum height=2.5cm, align=left, font=\tiny\sffamily, text width=2.5cm] (left) at (0,0) {
    This text is left\\aligned. Lines\\start from the\\left side.
  };
  \node[below=0.2cm of left, font=\scriptsize\sffamily\bfseries, text=ModuleBlue] {Left \key{Ctrl}+\key{L}};
  
  % Centre aligned
  \node[draw=FunPurple, fill=FunPurple!5, rounded corners=4pt, minimum width=3cm, minimum height=2.5cm, align=center, font=\tiny\sffamily, text width=2.5cm] (centre) at (4,0) {
    This text is\\centred. Every\\line is in the\\middle.
  };
  \node[below=0.2cm of centre, font=\scriptsize\sffamily\bfseries, text=FunPurple] {Centre \key{Ctrl}+\key{E}};
  
  % Right aligned
  \node[draw=FunCoral, fill=FunCoral!5, rounded corners=4pt, minimum width=3cm, minimum height=2.5cm, align=right, font=\tiny\sffamily, text width=2.5cm] (right) at (8,0) {
    This text is right\\aligned. Lines\\end at the\\right side.
  };
  \node[below=0.2cm of right, font=\scriptsize\sffamily\bfseries, text=FunCoral] {Right \key{Ctrl}+\key{R}};
  
  % Justify
  \node[draw=LabGreen, fill=LabGreen!5, rounded corners=4pt, minimum width=3cm, minimum height=2.5cm, align=justify, font=\tiny\sffamily, text width=2.5cm] (justify) at (12,0) {
    Justified text is\\flush on both the\\left and right\\sides evenly.
  };
  \node[below=0.2cm of justify, font=\scriptsize\sffamily\bfseries, text=LabGreen] {Justify \key{Ctrl}+\key{J}};
\end{tikzpicture}
\caption{Four text alignment options in MS Word}
\label{fig:alignment}
\end{figure}

\section{Line Spacing}
\label{sec:line-spacing}
\index{MS Word!line spacing}

% --- Table 25.1: Line Spacing Options ---
\begin{table}[H]
\centering
\caption{Line spacing options and when to use them}
\label{tab:line-spacing}
\renewcommand{\arraystretch}{1.3}
\begin{tabularx}{\textwidth}{l l X}
\toprule
\rowcolor{HeaderRow}
\textcolor{white}{\textbf{Spacing}} & \textcolor{white}{\textbf{Value}} & \textcolor{white}{\textbf{When To Use}} \\
\midrule
\rowcolor{RowLight} Single & 1.0 & Compact notes, short documents \\
\rowcolor{RowDark} 1.5 Lines & 1.5 & Most school assignments (recommended) \\
\rowcolor{RowLight} Double & 2.0 & Essays the teacher will annotate \\
\bottomrule
\end{tabularx}
\end{table}

\section{Lists}
\label{sec:lists}
\index{MS Word!lists}

Word supports two types of lists:
\begin{itemize}[leftmargin=*, label={\textcolor{ModuleBlue}{\faAngleRight}}, itemsep=2pt]
  \item \textbf{Bulleted lists} --- for items with no specific order (like this one!)
  \item \textbf{Numbered lists} --- for steps or items in a specific order
\end{itemize}

To create a list: Home tab $\rightarrow$ click the Bullets \faBars\ or Numbering \faListOl\ button.

\bytetip[fun]{The word ``font'' comes from the French word ``fonte,'' meaning ``something that has been melted.'' This is because early printing used melted metal to create letter shapes. Now we just click a dropdown menu --- much less dangerous!}

\begin{tryit}
Open a new Word document and type a title ``My Favourite Things.'' Make the title \textbf{bold}, \textit{size 18}, and \textcolor{ModuleBlue}{blue}. Below it, create a bulleted list of 5 things you love. Then change the alignment of the title to \textbf{centre}. Finally, set the line spacing to 1.5.
\end{tryit}

\begin{funfact}
There are over \textbf{200,000 fonts} available in the world! Windows comes with about 200 pre-installed. The most commonly used fonts for school work are \textbf{Calibri}, \textbf{Arial}, and \textbf{Times New Roman}.
\end{funfact}

\begin{reviewqs}
  \item Name four types of font formatting you can apply in Word.
  \item What are the four text alignment options and their shortcuts?
  \item What line spacing is recommended for school assignments?
  \item How do you create a bulleted list in Word?
  \item What is the shortcut for bold?
\end{reviewqs}

\begin{challenge}
Create a \textbf{one-page formatted document} about your school. It must include: a centred bold title (size 20), a subtitle in italic, two paragraphs with justified alignment, a bulleted list of school subjects, and at least 3 different font colours. Save it as ``MySchool.docx.''
\end{challenge}
